\begin{abstract}
3D Gaussian Splatting (3DGS) represents scenes through primitives with coupled intrinsic properties: geometric attributes (position, covariance, opacity) and appearance attributes (view-dependent color). 
Faithful reconstruction requires \textbf{intrinsic geometry-appearance consistency}, where geometry accurately captures 3D structure while appearance reflects photometry. 
However, sparse observations lead to appearance overfitting and underconstrained geometry, causing severe novel-view artifacts.
We present ICO-GS (Intrinsic Geometry-Appearance Consistency Optimization for 3DGS), a principled framework that enforces this consistency through tightly coupled geometric regularization and appearance learning. 
Our approach first regularizes geometry via feature-based multi-view photometric 
constraints by employing pixel-wise top-$k$ selection to handle occlusions and 
edge-aware smoothness to preserve sharp structures.
Then appearance is coupled with geometry through cycle-consistency depth filtering, which identifies reliable regions to synthesize virtual views that propagate geometric correctness into appearance optimization. 
Experiments on LLFF, DTU, and Blender show ICO-GS substantially improves geometry 
and photometry, consistently outperforming existing sparse-view baselines, particularly in challenging weakly-textured regions.
\end{abstract}

% Experiments on LLFF, DTU, and Blender show ICO-GS substantially improves geometry 
% and photometry, outperforming sparse-view baselines, especially in weakly-textured 
% regions.

% 3D Gaussian Splatting (3DGS) represents scenes through primitives with coupled intrinsic properties: geometric attributes (position, covariance, opacity) and appearance attributes (view-dependent color). Faithful reconstruction requires intrinsic geometry-appearance consistency, where geometry accurately captures 3D structure while appearance reflects photometry. However, sparse observations lead to appearance overfitting and underconstrained geometry, causing severe novel-view artifacts. We present ICO-GS (Intrinsic Geometry-Appearance Consistency Optimization for 3DGS), a principled framework that enforces this consistency through tightly coupled geometric regularization and appearance learning. Our approach first stabilizes geometry via feature-based multi-view photometric constraints—employing pixel-wise top-k selection to handle occlusions and edge-aware smoothness to preserve sharp structures. We then couple geometry with appearance through cycle-consistency depth filtering, which identifies reliable regions to synthesize virtual views that propagate geometric correctness into appearance optimization. Extensive experiments on LLFF, DTU, and Blender datasets demonstrate that ICO-GS substantially improves both multi-view geometry and novel-view photometry, consistently outperforming existing sparse-view baselines, particularly in challenging weakly-textured regions.



% 3D Gaussian Splatting (3DGS) represents scenes through primitives with coupled intrinsic properties: geometric attributes (position, covariance, opacity) and appearance attributes (view-dependent color). Faithful reconstruction requires intrinsic geometry-appearance consistency, where geometry accurately captures 3D structure while appearance reflects photometry. However, sparse observations lead to appearance overfitting and underconstrained geometry, causing severe novel-view artifacts.We present ICO-GS (Intrinsic Geometry-Appearance Consistency Optimization for 3DGS), a principled framework that enforces this consistency through tightly coupled geometric regularization and appearance learning. Our approach first regularizes geometry via feature-based multi-view photometric constraints by employing pixel-wise top-k selection to handle occlusions and edge-aware smoothness to preserve sharp structures.Then appearance is coupled with geometry through cycle-consistency depth filtering, which identifies reliable regions to synthesize virtual views that propagate geometric correctness into appearance optimization. Experiments on LLFF, DTU, and Blender show ICO-GS substantially improves geometry and photometry, consistently outperforming existing sparse-view baselines, particularly in challenging weakly-textured regions.